\documentclass[11pt]{article}
\usepackage{amsmath,amssymb,amsthm,mathtools}
\usepackage[margin=1in]{geometry}
\newtheorem{theorem}{Theorem}
\newtheorem{lemma}{Lemma}
\newtheorem{proposition}{Proposition}
\newtheorem{definition}{Definition}
\newtheorem{warning}{Warning}
\newtheorem{corollary}{Corollary}
\title{Step 84: CND Audit of the Paired Completed Weil Candidate}
\author{Six Birds / Formed-Layer Membrane Program}
\date{}
\begin{document}
\maketitle

\section{Purpose}
Steps 77--83 showed that separated diagonal accounting is the wrong target for the
RH membrane construction.  The prime and gamma slots should be kept as paired
shift-difference features whenever possible.  This step records the exact audit gate
for that strategy.

The active question is not merely whether a Fourier symbol is nonnegative.  The
log-side paired feature
\[
        q_\mu[f]=\int_{\mathbb R}\|\tau_a f-f\|_2^2\,d\mu(a)
\]
exists with a positive measure \(\mu\) exactly when its symbol is of
continuous-negative-definite/L\'evy--Khintchine type.  This is a stricter structural
condition than pointwise positivity.

\section{The CND paired-feature gate}
Let \(\tau_a f(t)=f(t+a)\).  If \(\mu\) is a positive symmetric measure satisfying the
usual L\'evy integrability condition, define
\[
        q_\mu[f]=\int_{\mathbb R}\|\tau_a f-f\|_2^2\,d\mu(a).
\]
On the Fourier side,
\[
        q_\mu[f]=\int_{\mathbb R}\Psi_\mu(\xi)|\widehat f(\xi)|^2\,d\xi,
        \qquad
        \Psi_\mu(\xi)=2\int_{\mathbb R}(1-\cos(a\xi))\,d\mu(a).
\]

\begin{definition}[Paired shift-feature admissibility]
A translation-invariant quadratic symbol \(\Psi\) is paired-shift admissible if
\[
        \Psi(\xi)=c\xi^2+2\int_{\mathbb R}(1-\cos(a\xi))\,d\mu(a)
\]
with \(c\ge0\), \(\mu\ge0\), and the standard L\'evy integrability condition.
Equivalently, \(\Psi\) is an even continuous negative definite symbol with
\(\Psi(0)=0\).
\end{definition}

\begin{theorem}[CND gate]
A translation-invariant positive paired feature of the above log-shift form exists
for \(\Psi\) if and only if \(\Psi\) is paired-shift admissible.  Equivalently,
\(e^{-t\Psi}\) is positive definite for every \(t>0\).
\end{theorem}

\begin{proof}
This is the symmetric one-dimensional L\'evy--Khintchine/Schoenberg theorem.  If
\(\Psi\) has the displayed representation, then
\(e^{-t\Psi}\) is positive definite and the quadratic form is realized by the
corresponding shift-difference feature.  Conversely, a continuous negative definite
even function with \(\Psi(0)=0\) has the L\'evy--Khintchine representation and hence
is realized by the paired feature.
\end{proof}

\section{Prime and gamma slots}
For a finite prime-power cutoff, set
\[
        \Psi_{{\rm pr},L}(\xi)
        =2\sum_n w_{n,L}(1-\cos(\xi\log n)),
        \qquad
        w_{n,L}\ge0.
\]
This is paired-shift admissible with the discrete measure
\(\sum_n w_{n,L}\delta_{\log n}\) and its symmetric reflection.

For the zeta gamma factor one obtains the normalized positive symbol
\[
        \Psi_\Gamma(\xi)
        =\frac12\int_0^\infty \frac{e^{-t/4}}{1-e^{-t}}
          \left(1-\cos\frac{\xi t}{2}\right)\,dt.
\]
It is also paired-shift admissible, represented by the positive kernel
\[
        d\mu_\Gamma(a)
        =\frac12\frac{e^{-2a/4}}{1-e^{-2a}}\,2\,da
\]
under the change \(a=t/2\), with the paired-difference form keeping the near-zero
singularity lawful.

\begin{proposition}[Prime--gamma CND closure]
For each finite prime cutoff \(L\),
\[
        \Psi_{{\rm pg},L}:=\Psi_{{\rm pr},L}+\Psi_\Gamma
\]
is paired-shift admissible.  Therefore the prime-plus-gamma paired carrier admits a
positive shift-feature representation.
\end{proposition}

\begin{proof}
The cone of continuous negative definite symbols is closed under nonnegative linear
combination.  Both \(\Psi_{{\rm pr},L}\) and \(\Psi_\Gamma\) are in that cone.
\end{proof}

\section{What this does not earn}
The CND audit only proves that a candidate prime-plus-gamma paired carrier is a
lawful positive feature.  It does not prove that this carrier is equal to the signed
completed Weil explicit-formula form.  The matching gate remains
\[
        q_Z[f]\le q_{{\rm pr}+\Gamma,L}[f]+q_{\rm pole}[f]+q_{\rm tail}[f]+e_L[f]
\]
on a completed anti-invariant core, followed by closability, fixed-ledger, and tail
records.

\begin{warning}[Nonnegative is not enough]
A pointwise nonnegative symbol need not be paired-shift admissible.  For example,
\(\Psi(\xi)=\xi^4\) is nonnegative but not continuous negative definite.  With
points \(-1,0,1\) and coefficients \((1,-2,1)\), whose sum is zero,
\[
        \sum_{i,j}c_i c_j \Psi(x_i-x_j)=24>0,
\]
violating the defining CND inequality.
\end{warning}

\begin{warning}[Finite-rank repair cannot fix a full-core non-CND residual]
A finite-rank pole/completion form can repair finitely many endpoint or null-mode
directions, but it cannot convert a non-CND infinite-dimensional residual into a
paired shift-feature on the full completed core.  Such a residual must be repaired
by a full-sector tail/source, a different paired representation, or recorded as a
defect/nonclaim.
\end{warning}

\section{RH gate status after this step}
The candidate paired feature package now has the following status.
\begin{itemize}
\item The finite prime slot passes the paired-feature gate.
\item The normalized gamma slot passes the paired-feature gate as a paired, singular,
renormalized difference form.
\item Their sum is CND and therefore has a positive paired feature representation.
\item The pole slot remains finite-rank and is lawful only for finite completion/null
modes.
\item The tail/support slot must supply full-sector coercivity or a vanishing defect
record.
\item The signed-to-positive equality with the actual completed Weil form remains
unearned.
\end{itemize}

\section{Conclusion}
The active RH construction target is now sharpened:
\[
        \boxed{\text{prove that the actual completed Weil form equals or is dominated by
        a paired CND carrier plus lawful pole/tail records}.}
\]
The CND audit says the prime-plus-gamma candidate is a legitimate positive feature,
but not yet the accepted Douglas bridge.

\end{document}
